\documentclass[11pt]{article}

\usepackage[margin=1in]{geometry}
\usepackage{graphicx}
\usepackage{booktabs}
\usepackage{amsmath}
\usepackage{xcolor}
\usepackage{caption}
\usepackage{float}
\usepackage{hyperref}
\usepackage{authblk}
\hypersetup{colorlinks=true, linkcolor=blue, citecolor=blue, urlcolor=blue}

\title{PE-RankFormer: A Minimally-Mechanistic Relational Transformer\\
for Prime-Editing Efficiency Prediction}
\author{PEFormer research pilot}
\date{\today}

\begin{document}
\maketitle

\begin{abstract}
Prime editing (PE) efficiency prediction has been dominated by mechanistic models
that hard-code assumptions about the PE reaction pathway (nick, flap resolution,
mismatch repair, etc.), most recently OptiPrime (Hsu et al., 2026), a strongly
mechanistic model trained on 297{,}962 PE experiments. We ask whether a
minimally-mechanistic, general-purpose relational Transformer, given the same
scale of training data, can match or exceed it. We reconstruct 262{,}508 of the
297{,}962 target training examples (88.1\%) from four source studies, implement
PE-RankFormer -- a 26.6M-parameter paired-sequence Transformer with cross-attention
between the edit-outcome encoding and the pegRNA encoding and FiLM-based
conditioning on experimental context -- and train it with a joint regression and
within-target ranking objective. A single architectural addition motivated by
diagnosing early overfitting, a 3-way categorical ``simplex'' outcome head that
models the mutually-exclusive unedited/edited/indel read fractions instead of
regressing edited-read fraction alone, produced the largest single gain in the
project. A 6-model ensemble of PE-RankFormer variants, evaluated head-to-head
against OptiPrime's own released 5-model ensemble on the 15{,}022-row Hsu subset
of the locked test fold, achieves Spearman $\rho = 0.807$ versus OptiPrime's
$0.724$ (Pearson $0.780$ vs.\ $0.724$), a difference of $+0.083$ with a
protospacer-clustered paired bootstrap 95\% CI of $[+0.047, +0.118]$
($p < 0.0005$, PE-RankFormer wins in 2000/2000 resamples). We report this result
together with negative results from the same sweep and explicit caveats about
corpus completeness and test-fold query discipline.
\end{abstract}

\tableofcontents

% ============================================================
\section{Introduction}
% ============================================================

Prime editing efficiency prediction is typically approached with strongly
mechanistic models: pipelines that explicitly encode the PE reaction as a
sequence of biochemical steps (spacer/PBS/RTT annealing, flap competition,
mismatch-repair outcome) and use that structure to constrain the model. OptiPrime
(Hsu et al., 2026) is the current state of the art on this axis, trained on
297{,}962 PE experiments assembled across four studies.

The research question motivating this pilot is whether a \emph{minimally
mechanistic} model -- one that encodes only the raw sequences and experimental
metadata and lets a general-purpose relational architecture (a Transformer with
cross-attention between the edit and pegRNA representations) discover whatever
structure is predictive -- can match or outperform a strongly mechanistic model
of comparable training scale. This is a real test of an increasingly common
hypothesis in computational biology: that architectural inductive bias from
domain mechanism can be substituted for by scale plus a sufficiently expressive
general-purpose architecture, provided the objective is designed to exploit the
structure actually present in the data (e.g., within-target ranking, or -- as
turned out to matter most in this pilot -- the full space of mutually exclusive
sequencing outcomes rather than a single scalar summary of them).

This report documents the full pipeline: corpus reconstruction, architecture,
training procedure, a systematic ablation and configuration sweep, and a
statistically validated, ensemble-matched head-to-head comparison against
OptiPrime's own released baseline.

% ============================================================
\section{Training Corpus Reconstruction}
% ============================================================

OptiPrime was trained on 297{,}962 PE experiments assembled from four source
studies. We reconstructed this corpus independently from public data releases,
verified against a 42-partition ground-truth breakdown (by study $\times$ cell
type $\times$ PE system) reported in the OptiPrime paper.

\begin{table}[H]
\centering
\caption{Final reconstructed corpus composition
(\texttt{data/processed/optiprime\_full\_297962.parquet}).}
\begin{tabular}{lrr}
\toprule
Source & Partitions matched (of target) & Rows \\
\midrule
Liu (Hsu 2026, newly generated) & 4/4 (raw count; see note) & 74{,}769 \\
Kim (DeepPrime) & 18/18 exact & 58{,}301 \\
Schwank (PRIDICT / PRIDICT2.0) & 18/20 exact & 129{,}438 \\
\midrule
\textbf{Total} & \textbf{36/42 target-exact} & \textbf{262{,}508} \\
\bottomrule
\end{tabular}
\end{table}

262{,}508 of 297{,}962 target rows (88.1\%) were recovered with a documented,
reproducible provenance trail. The three remaining discrepancies were
investigated and ruled out rather than papered over: (1) the Hsu/Liu partition
contains 74{,}769 raw rows against a 65{,}594-row target with no public QC/coverage
column to explain the gap -- several plausible filters (PE2/PE4 co-presence,
barcode-collision exclusion, invalid-read exclusion, zero-efficiency exclusion,
OptiPrime's own embedding-size bounds) were each tested and rejected, so we use
the full 74,769; (2) the Schwank K562 PE4+epegRNA partition (target 21{,}201) is
essentially absent from any public release (823 rows found, after checking
PRIDICT v1's full supplementary branch, PRIDICT2.0's processed data, and the
ranking-percentile file); (3) the Schwank K562 PE2+epegRNA partition has a small
94-row excess traced to an unapplied PRIDICT2.0 QC exclusion list. Full detail
in \texttt{reports/training\_data\_reconstruction.md} and
\texttt{reports/dataset\_reconstruction\_status.md}.

Cross-validation folds were built to be \textbf{protospacer-disjoint}: fold 0
(test), fold 1 (val), folds 2--4 (train), seed 20260812, with zero-leakage
verified (no protospacer appears in more than one fold).

% ============================================================
\section{Model Architecture}
% ============================================================

PE-RankFormer is a 26.6M-parameter Transformer ($d_{\text{model}}{=}384$,
6 attention heads, FFN width 1536, dropout 0.10) with two parallel token
encoders and a bidirectional cross-attention stage:

\begin{itemize}
\item \textbf{Edit encoder} (6 layers): a paired-sequence encoding of the
  unedited and edited full target sequence, so the model can directly attend
  to what changed rather than being told a structured description of the edit.
\item \textbf{PegRNA encoder} (4 layers): a segment-aware encoding of the
  spacer, PBS, and RTT, with segment-type embeddings distinguishing the three
  regions.
\item \textbf{Cross-attention} (2 blocks): bidirectional cross-attention between
  the edit-token and pegRNA-token representations, followed by attention
  pooling on each side.
\item \textbf{FiLM conditioning}: pooled representations are modulated by an
  embedding of experimental context (cell type, PE system variant, Cas9/PAM
  variant, scaffold, motif), so the same sequence pair can receive different
  predictions in different experimental contexts.
\item \textbf{Prediction head}: either a scalar sigmoid head (edited-fraction
  regression) or the 3-way simplex head described in \S\ref{sec:simplex}.
\end{itemize}

No PE-reaction-specific mechanism (e.g., explicit flap-competition or
mismatch-repair modeling) is encoded; all such structure, if present in the
data, must be discovered by the cross-attention mechanism itself -- this is
the ``minimally mechanistic'' design constraint against which the model is
being tested.

\subsection{The simplex outcome head: diagnosis-driven innovation}
\label{sec:simplex}

Early training runs (Model B, scalar head, no ranking loss) showed clear
overfitting: training loss fell roughly $4\times$ over 30 epochs while
validation Spearman $\rho$ plateaued from epoch $\sim$20 onward
(Figure~\ref{fig:training-curves}). Adding model capacity would not address
this; the constraint was in the supervision signal, not the architecture's
expressive power.

Every PE sequencing read falls into exactly one of three mutually exclusive
categories: \emph{unedited}, \emph{edited} (the intended outcome), or
\emph{indel} (an unintended byproduct of the nick/repair process). The original
scalar-head design discarded the indel read fraction entirely and regressed
only the edited fraction -- despite indel being present with 100\% coverage in
the reconstructed corpus, nonzero in $\sim$60\% of rows, and correlated with
editing efficiency ($r{\approx}0.25$). This is a real, cheaply available signal
being thrown away.

We replaced the scalar sigmoid head with a \textbf{3-way simplex head}: a
softmax over (unedited, edited, indel) logits trained with soft cross-entropy
against the observed read proportions (clipped/renormalized on the rare rows
where edited+indel exceeds 1 due to counting noise). At inference,
$\text{efficiency} = \text{softmax}(\text{logits})_{\text{edited}}$, and the
ranking score used for the pairwise ranking loss is the edited-vs-unedited
log-odds, $\text{logit}_{\text{edited}} - \text{logit}_{\text{unedited}}$.

This does not violate the minimally-mechanistic constraint: mutual exclusivity
of the three read categories is a property of how sequencing reads are
classified (an assay-level fact), not an assumption about the underlying PE
reaction graph. No biochemical step ordering or intermediate state is encoded.

The simplex head produced the top 4 single-model validation Spearman results
of the entire sweep (\S\ref{sec:sweep}). Two further variations, tried to see
whether the gain could be pushed further, produced negative results and are
reported as such rather than discarded: training the regression component in
clipped-logit space instead of raw $[0,1]$ space gave no improvement
(0.866 vs.\ 0.864 raw, within seed noise), and increasing dropout to 0.20
gave no improvement (0.871 vs.\ 0.870 at dropout 0.10).

% ============================================================
\section{Training Procedure}
% ============================================================

Models are trained with joint regression (Huber loss) and pairwise RankNet
within-target ranking loss ($\lambda_{\text{rank}}$ weighting), using a custom
grouped batch sampler that biases minibatches toward rows sharing a
within-target ranking group -- necessary because the $\sim$220{,}000 distinct
ranking groups in the full corpus mean uniform random batching yields close to
zero rankable pairs per batch. Optimization uses AdamW with a warmup + cosine
schedule, BF16 mixed precision, and fused scaled-dot-product attention, on a
single RTX PRO 6000 Blackwell GPU. All models in the final sweep were trained
for a full 30 epochs (early-stopping patience matched across configurations
after an initial unfair comparison was caught and corrected -- see
\S\ref{sec:caveats}).

% ============================================================
\section{Ablation and Configuration Sweep}
\label{sec:sweep}
% ============================================================

A bounded sweep over outcome head (scalar vs.\ simplex), regression target
space (raw $[0,1]$ vs.\ clipped-logit), ranking-loss weight $\lambda_{\text{rank}}$,
context conditioning (on/off), and dropout was run, each configuration trained
to full 30 epochs at matched patience. All selection decisions in this sweep
were made on the \textbf{validation} fold only.

\begin{table}[H]
\centering
\caption{Full configuration sweep, sorted by best validation Spearman $\rho$
(single seed 20260812 unless noted). ``rank'' = with ranking loss
($\lambda_{\text{rank}}{=}0.25$ for Models A/C, $0.05$ for the low-rank probe);
blank = ranking loss off ($\lambda_{\text{rank}}{=}0$).}
\label{tab:sweep}
\begin{tabular}{lllrrr}
\toprule
Run & Head & Space & Context & $\lambda_{\text{rank}}$ & Val $\rho$ \\
\midrule
final\_simplex\_s2      & simplex & raw    & on  & 0.00 & \textbf{0.879} \\
final\_simplex\_s3      & simplex & raw    & on  & 0.00 & 0.873 \\
exp\_dropout20\_simplex & simplex & raw    & on  & 0.00 & 0.871 (dropout 0.20) \\
exp\_simplex            & simplex & raw    & on  & 0.00 & 0.870 \\
model\_b\_norank        & scalar  & raw    & on  & 0.00 & 0.864 \\
final\_logit\_s3        & scalar  & logit  & on  & 0.00 & 0.867 \\
exp\_logit\_norank      & scalar  & logit  & on  & 0.00 & 0.866 \\
final\_logit\_s2        & scalar  & logit  & on  & 0.00 & 0.862 \\
exp\_lowrank005         & scalar  & raw    & on  & 0.05 & 0.850 \\
model\_a\_rank          & scalar  & raw    & on  & 0.25 & 0.786 \\
model\_c\_nocontext\_fair & scalar & raw  & off & 0.25 & 0.705 \\
model\_c\_nocontext     & scalar  & raw    & off & 0.25 & 0.675 (early-stopped, unfair) \\
\bottomrule
\end{tabular}
\end{table}

Three findings stand out. First, the simplex head is the single largest lever
in the sweep, outperforming every scalar-head configuration regardless of
other settings. Second, the ranking loss \emph{hurts} single-model validation
Spearman at the weights tried ($\lambda_{\text{rank}}{=}0.25$ costs $\sim$0.08
$\rho$ relative to $\lambda_{\text{rank}}{=}0$); this is a genuine trade-off
against within-target ranking quality that a scalar global Spearman number does
not fully capture (see \S\ref{sec:full-test} for within-target metrics), and is
noted as a caveat rather than resolved in this pilot. Third, removing
experimental-context conditioning is unambiguously costly (0.705 vs.\ 0.786 at
matched patience, both with ranking loss on), confirming the FiLM conditioning
mechanism is doing real work.

\begin{figure}[H]
\centering
\includegraphics[width=0.85\textwidth]{../results/figures/06_training_curves.pdf}
\caption{Training loss and validation Spearman $\rho$ for the three initial
scalar-head configurations: Model A (rank loss on), Model B (rank loss off),
Model C (no context conditioning). Model B's training loss falls
$\sim\!4\times$ while validation performance plateaus after epoch $\sim$20 --
the overfitting signal that motivated the simplex-head redesign
(\S\ref{sec:simplex}).}
\label{fig:training-curves}
\end{figure}

% ============================================================
\section{Headline Result: Ensemble vs.\ Ensemble Comparison with OptiPrime}
% ============================================================

OptiPrime's own reported baseline number is itself a 5-model ensemble. To make
a fair, scale-matched comparison rather than pitting a single PE-RankFormer
checkpoint against an ensemble, we built our own \textbf{6-model ensemble}
(3 simplex-head seeds + 3 logit-space scalar-head seeds, selected on validation
only) and evaluated it once, head-to-head against OptiPrime's real released
weights, on the 15{,}022-row Hsu-only subset of the locked test fold (the
subset on which OptiPrime's own feature cache and code could be run without
modification).

\begin{table}[H]
\centering
\caption{Global correlation on the locked Hsu test subset ($n=15{,}022$).}
\label{tab:headline}
\begin{tabular}{lrrrr}
\toprule
Model & Pearson $r$ & Spearman $\rho$ & MAE & RMSE \\
\midrule
OptiPrime (official 5-model ensemble) & 0.724 & 0.724 & 0.091 & 0.129 \\
\textbf{PE-RankFormer ENSEMBLE (6-model)} & \textbf{0.780} & \textbf{0.807} & \textbf{0.078} & \textbf{0.117} \\
PE-RankFormer (no-rank, single) & 0.756 & 0.775 & 0.083 & 0.125 \\
PE-RankFormer (rank, single) & 0.659 & 0.670 & 0.101 & 0.147 \\
\bottomrule
\end{tabular}
\end{table}

\begin{figure}[H]
\centering
\includegraphics[width=0.85\textwidth]{../results/figures/02_global_correlation_comparison.pdf}
\caption{Global correlation comparison on the locked Hsu test subset. The
PE-RankFormer ensemble (orange/blue) exceeds OptiPrime's official ensemble
(grey) on both Pearson and Spearman correlation; dotted line marks OptiPrime's
Spearman $\rho$ for reference.}
\label{fig:headline}
\end{figure}

\subsection{Statistical significance}

We ran a paired, protospacer-clustered bootstrap (2000 resamples, seed
20260812) on the identical 15{,}022 held-out rows, resampling whole
protospacers rather than rows -- the correct unit of statistical dependence,
since many pegRNA designs share a protospacer and a naive row-level or
ranking-group-level bootstrap would understate the true confidence interval
(in the Hsu library design, ranking groups are almost all singletons, so a
group-level bootstrap is effectively a disguised row-level one).

\begin{table}[H]
\centering
\caption{Paired target-level bootstrap, PE-RankFormer ENSEMBLE vs.\ OptiPrime.}
\begin{tabular}{lr}
\toprule
Paired rows & 15{,}022 \\
Protospacer clusters & 233 \\
PE-RankFormer Spearman $\rho$ & 0.8069 \\
OptiPrime Spearman $\rho$ & 0.7242 \\
Observed difference & $+0.0827$ \\
Bootstrap mean difference & $+0.0838$ \\
95\% CI & $[+0.0470,\ +0.1176]$ \\
Fraction of resamples PE-RankFormer wins & $1.0000$ (2000/2000) \\
Two-sided bootstrap $p$ & $< 0.0005$ \\
\bottomrule
\end{tabular}
\end{table}

The confidence interval excludes zero by a wide margin and PE-RankFormer wins
in every one of 2000 bootstrap resamples. This is a decisive, statistically
validated win over OptiPrime's own released baseline on its own evaluation
subset.

\subsection{Full test fold metrics}
\label{sec:full-test}

Beyond the Hsu-only OptiPrime-comparable subset, the ensemble was also
evaluated on the full locked test fold ($n=52{,}319$, spanning all four source
studies), where within-target ranking metrics are meaningful (the Hsu subset
alone has almost no multi-design-per-target groups).

\begin{table}[H]
\centering
\caption{PE-RankFormer ENSEMBLE, full locked test fold ($n=52{,}319$).}
\begin{tabular}{lr}
\toprule
Global Pearson $r$ & 0.875 \\
Global Spearman $\rho$ & 0.894 \\
MAE / RMSE & 0.082 / 0.131 \\
Within-target Spearman $\rho$ (macro mean) & 0.541 [0.496, 0.583] \\
Top-1 regret & 0.0113 [0.0101, 0.0127] \\
Top-3 regret & 0.00017 [0.00009, 0.00027] \\
Top-1 / Top-3 recall & 0.629 / 0.969 \\
NDCG@3 / NDCG@5 & 0.920 / 0.931 \\
\bottomrule
\end{tabular}
\end{table}

\begin{figure}[H]
\centering
\includegraphics[width=0.7\textwidth]{../results/figures/01_predicted_vs_observed.pdf}
\caption{Predicted vs.\ observed editing efficiency, full locked test fold.
Colorbar is log-scaled to avoid the low-efficiency mode dominating the visual
impression. The decile-mean calibration overlay (black line) tracks $y=x$
closely through $\sim$60\% observed efficiency, with under-prediction at the
high-efficiency tail.}
\label{fig:calibration}
\end{figure}

\begin{figure}[H]
\centering
\includegraphics[width=0.85\textwidth]{../results/figures/05_performance_by_context.pdf}
\caption{PE-RankFormer global Spearman $\rho$ broken down by experimental
context (cell type $\times$ PE system), full test fold, contexts with
$n\geq20$.}
\label{fig:by-context}
\end{figure}

\begin{figure}[H]
\centering
\includegraphics[width=0.85\textwidth]{../results/figures/03_within_target_spearman_distribution.pdf}
\caption{Distribution of within-target Spearman $\rho$ across target groups
with $n\geq5$ designs, full test fold.}
\label{fig:within-target}
\end{figure}

\begin{figure}[H]
\centering
\includegraphics[width=0.6\textwidth]{../results/figures/07_ablation_comparison.pdf}
\caption{Global Spearman, within-target Spearman, and top-3 regret across the
three single-model configurations and the final ensemble, full test fold. The
ensemble improves all three metrics simultaneously relative to any single
constituent model.}
\label{fig:ablation}
\end{figure}

% ============================================================
\section{Discussion and Caveats}
\label{sec:caveats}
% ============================================================

\textbf{Corpus completeness.} The reconstructed corpus covers 88.1\% of
OptiPrime's target training set (262{,}508 / 297{,}962 rows), not 100\%. The
missing 12\% is concentrated in two Hsu/Liu and Schwank partitions for which no
public data release appears to contain the missing rows (\S2), rather than
being randomly missing. This is disclosed as a limitation on the training-data
side of the comparison; both models were nonetheless trained and evaluated on
the same fold structure, and the head-to-head comparison itself uses
OptiPrime's own real released weights (not a reproduction), so the effect of
any corpus gap runs in OptiPrime's favor if anything (PE-RankFormer had
slightly less data than the original OptiPrime training run).

\textbf{Test-fold query discipline.} All final model \emph{selection} decisions
(which configuration, which ensemble composition) were made on the validation
fold only, and the final headline ensemble-vs-ensemble number was computed
once on the locked test fold. However, the test fold was queried multiple
times earlier in the project during the Model A/B/C ablation phase before the
validation-only selection discipline was adopted for the final sweep. This is
a mild multiple-comparisons exposure, disclosed rather than hidden; the
magnitude of the final effect ($+0.083$ Spearman, CI excluding zero by a wide
margin) is unlikely to be an artifact of it, but it is not a pre-registered
single test.

\textbf{Ranking-loss trade-off.} The ranking loss, as configured
($\lambda_{\text{rank}}\in\{0.05, 0.25\}$), costs global validation Spearman
relative to turning it off, and was accordingly not used in the final ensemble
(all 6 members trained with $\lambda_{\text{rank}}{=}0$). This does not mean
within-target ranking is unimportant -- the final ensemble's within-target
Spearman (0.541) and top-3 regret (0.00017) on the full test fold are strong --
but the specific pairwise RankNet formulation and weight tried here did not
improve on what the simplex head's soft-label supervision already provides.
Whether a better-tuned ranking loss (lower weight, different formulation, or
combined with the simplex head) could add further gains is an open question
for future work.

\textbf{Calibration at high efficiency.} Figure~\ref{fig:calibration} shows
good calibration through $\sim$60\% observed efficiency with systematic
under-prediction in the high-efficiency tail ($>$80\%). This is a known-shape
limitation to flag for any downstream use case that specifically needs
accurate absolute predictions in that regime (as opposed to relative ranking,
where top-3 regret is very low).

% ============================================================
\section{Conclusion}
% ============================================================

A minimally-mechanistic relational Transformer, trained on an independently
reconstructed 88.1\%-complete version of OptiPrime's training corpus, matches
and then clearly exceeds OptiPrime's own released 5-model ensemble on its own
evaluation subset: Spearman $\rho = 0.807$ vs.\ $0.724$
($+0.083$, 95\% CI $[+0.047, +0.118]$, $p<0.0005$, 2000/2000 bootstrap wins).
The single most impactful design decision was not scaling model capacity --
which early diagnosis showed would not help, given clear overfitting -- but
better exploiting the label structure already present in the data via a 3-way
simplex outcome head that models the full mutually-exclusive read-outcome
distribution instead of a single scalar summary of it. This supports the
pilot's central hypothesis: with careful supervision design and honest,
validation-disciplined model selection, a general-purpose relational
architecture without hard-coded PE-reaction mechanism can outperform a
strongly mechanistic model at comparable training scale.

\end{document}
